\section{Conclusion}
\label{sec:conclusion}

The preceding GFM papers define what to maximize ($\volL(G)$), prove
structural safety properties, and establish anti-monopolar pressure. This
paper addresses how the actor \emph{reasons about other agents}: providing
a causal model that upgrades scorpion detection and risk evaluation from
statistical to causal identification and sharpens the remaining
evaluation mechanisms.

Causal attribution (Section~\ref{sec:causal}) upgrades the foundational
paper's correlational scorpion detection to causal identification through
do-calculus counterfactuals. The convergence advantage over correlational
detection is proportional to confounding variance, which is substantial in
populations with many simultaneously acting agents. Non-stationary scorpion
bounds (Section~\ref{sec:scorpion}) characterize the evasion bandwidth: how
fast a scorpion must adapt to remain undetected, and the maximum undetected
contraction rate.

Risk-trust dynamics (Section~\ref{sec:risk_trust}) provide formal EWMA
update equations for $\Trisk_j$, with a two-gate architecture separating
attention allocation (log-odds aggregation) from action decisions (the
actor's own causal verification). This prevents coordinated report-bombing from being amplified into
action while preserving sensitivity to genuine independent agreement. Probabilistic
dependency risk (Section~\ref{sec:dependency}) converts minimax substrate
bounds into per-step expected costs that compound at the same rate as
growth, providing the quantitative risk model needed for risk-adjusted
planning.

Together with the foundational framework \citep{lasser2026gfm}, the
computable poset \citep{lasser2026poset}, and the horizon-aware objective
\citep{lasser2026horizon}, these results narrow the gap between the GFM
framework's structural guarantees and the mechanisms a concrete agent
would need to realize them. The results here are conditional on seven load-bearing assumptions:
causal sufficiency, additive SCM structure, and exogenous
orthogonality (Proposition~\ref{prop:causal_dominates}(c));
conditional independence of reporters given a shared prior
(Definition~\ref{def:logodds});
cascade confinement within substrates and the adversarial balance
condition (Proposition~\ref{prop:expected_risk}); and a
stationary-epoch adversary model
(Proposition~\ref{prop:evasion_bandwidth}). Each
assumption is explicit, falsifiable, and load-bearing.
As implementation machinery, these results do not affect the
framework's structural guarantees: a flaw in any mechanism here
degrades the actor's evaluation accuracy but does not compromise
the self-balancing, anti-monopolar, or structural alignment
properties established in Papers~1--3
\citep{lasser2026gfm, lasser2026poset, lasser2026horizon}.
The remaining
work is both theoretical (relaxing each assumption and characterizing
what happens when it fails) and empirical: instantiating these
computations in a concrete environment and checking whether the
predicted alignment properties hold under real conditions.
